% Figure: merit-order supply curve (dispatch stack as a step function of
% cumulative capacity versus marginal cost) crossed by a vertical demand
% line. The intersection sets the system marginal price. Step plot built
% with const plot. Values illustrative, current as of 2025.
\begin{figure}[htbp]
\centering
\begin{tikzpicture}
\begin{axis}[
  width=13cm, height=8cm,
  xlabel={cumulative available capacity (GW)},
  ylabel={marginal cost (USD/MWh)},
  xmin=0, xmax=50, ymin=0, ymax=200,
  xtick={0,10,20,30,40,50},
  ytick={0,40,80,120,160,200},
  axis lines=left,
  tick label style={font=\scriptsize},
  label style={font=\small},
  legend style={font=\scriptsize, at={(0.03,0.97)}, anchor=north west, draw=none},
]
% merit-order step curve: cheap renewables and nuclear first, gas last
\addplot[engblue, very thick, const plot] coordinates {
  (0,0)(12,0)(12,8)(20,8)(20,35)(30,35)(30,70)(40,70)(40,140)(48,140)(48,200)(50,200)
};
\addlegendentry{merit-order supply}
% demand line at 34 GW
\addplot[engred, very thick, dashed] coordinates {
  (34,0)(34,200)
};
\addlegendentry{demand $= 34$ GW}
% marker at the clearing point
\addplot[engorange, only marks, mark=*, mark size=2pt] coordinates {(34,70)};
\node[anchor=west, font=\scriptsize, color=engorange]
  at (axis cs:34.5,82) {clearing price $\approx 70$};
\end{axis}
\end{tikzpicture}
\caption[Merit-order supply curve and marginal price]{The merit-order
supply curve for an illustrative grid: generators stacked left to right in
order of rising marginal cost, from near-zero-cost wind and solar through
nuclear and coal to flexible gas (current as of 2025). The vertical demand
line cuts the stack at the marginal generator, and its marginal cost sets
the system clearing price that every dispatched generator is paid. As
demand rises into the steep right of the stack, the price climbs sharply,
which is the mechanism behind scarcity pricing and the case for demand
response. Quantities illustrative.}
\label{fig:vol04ch14:merit-order}
\end{figure}
